\section{Проектирование и архитектура программного комплекса}

В первом разделе рассмотрены теоретические основы сбора, подготовки и визуализации данных о недвижимости. Раздел посвящён проектным решениям: требованиям к системе, составу компонентов, хранению данных, серверному интерфейсу, ролям пользователей и экранам дашборда.

\subsection{Требования к программному комплексу}

Исходя из постановки задачи (см.~введение), сформулированы функциональные и нефункциональные требования.

\textbf{Функциональные требования:}
\begin{enumerate}
    \item автоматизированный сбор объявлений с платформы Яндекс.Недвижимость с возможностью поэтапного выполнения и возобновления после остановки или капчи;
    \item подготовка табличного датасета с единой схемой полей (цена, площадь, комнатность, координаты, параметры дома, категориальные признаки);
    \item хранение серверных снимков рынка по нескольким городам и загрузка пользовательских CSV-файлов;
    \item интерактивный дашборд: фильтры, сводные показатели, настраиваемые графики, карта объектов, сравнение двух сегментов рынка, таблица объявлений;
    \item экспорт отфильтрованных данных и графиков;
    \item обратное геокодирование координат в адрес с кэшированием результатов;
    \item разграничение доступа по ролям: просмотр, расширенная аналитика, администрирование данных и парсера, управление пользователями;
    \item удалённый запуск стадий парсера и контроль прогресса из веб-интерфейса администратора.
\end{enumerate}

\textbf{Нефункциональные требования:}
\begin{itemize}
    \item \emph{отзывчивость интерфейса} --- смена фильтров не должна блокировать ввод; пересчёт показателей выполняется на клиенте без полной перезагрузки страницы;
    \item \emph{согласованность фильтрации} --- одни и те же правила отбора строк применяются в интерфейсе, при серверной пагинации и экспорте;
    \item \emph{масштаб выборки} --- поддержка десятков тысяч объявлений в одном городском снимке (до 50~000 строк);
    \item \emph{развёртывание без отдельного СУБД-сервера} --- для учебного и локального использования достаточно встроенной SQLite и файлового хранения таблиц;
    \item \emph{расширяемость по городам} --- добавление нового рынка сводится к конфигурации парсера и каталогу данных без переписывания клиента.
\end{itemize}

На момент проектирования в систему включены два города --- Москва и Саратов; архитектура рассчитана на добавление других регионов по тому же шаблону.

\subsection{Общая архитектура системы}

Программный комплекс построен по схеме «клиент --- сервер --- внешний источник» с выделенным автономным модулем сбора данных. Логически выделяются четыре части:
\begin{itemize}
    \item \textbf{клиентское веб-приложение} --- одностраничное приложение (SPA) в браузере: авторизация, дашборд, панели администратора;
    \item \textbf{сервер приложений} --- REST API на Node.js: выдача датасетов, авторизация, геокодирование, управление задачами парсера и учётными записями;
    \item \textbf{модуль парсера} --- Python-пайплайн на базе Selenium: обход сайта, извлечение карточек, ETL и контроль качества;
    \item \textbf{общий слой типов и правил} --- каталог с общими для клиента и сервера описаниями городов, фильтров и форматов данных, чтобы логика не расходилась между уровнями.
\end{itemize}

Схема взаимодействия компонентов представлена на рисунке~\ref{fig:architecture}.

\begin{figure}[htbp]
\centering
\begin{tikzpicture}[
    comp/.style={draw, rounded corners=2pt, align=center, minimum height=12mm, inner sep=5pt, font=\small},
    store/.style={comp, fill=gray!10, minimum width=32mm, text width=30mm},
    proc/.style={comp, fill=blue!8, minimum width=34mm, text width=32mm},
    ext/.style={comp, fill=orange!12, dashed, minimum width=34mm, text width=32mm},
    lbl/.style={font=\scriptsize, fill=white, inner sep=1.5pt},
    arr/.style={-{Stealth[length=2.2mm]}, semithick}
]
\node[proc] (client) {Браузер\\React SPA};
\node[proc, below=9mm of client] (api) {Express API\\порт 3001};

\coordinate (rowcenter) at ($(api.south) + (0,-14mm)$);
\node[store] (json) at ([xshift=-42mm]rowcenter) {JSON\\датасеты\\по городам};
\node[store] (sqlite) at (rowcenter) {SQLite\\users, jobs, cache};
\node[proc] (parser) at ([xshift=42mm]rowcenter) {Python-парсер\\Selenium + pandas};

\node[fit=(json)(sqlite)(parser), inner sep=6pt, draw=none] (bottomrow) {};

\node[ext, below=12mm of parser] (yandex) {realty.yandex.ru};
\node[ext, right=24mm of api] (geocoder) {Геосервисы\\Яндекс, OSM};

\draw[arr] (client) -- node[right, lbl] {HTTP /api} (api);
\coordinate (apibranch) at ($(api.south)!0.35!(rowcenter)$);
\draw[arr] (api.south) -- (apibranch);
\draw[arr] (apibranch) -| (json.north);
\draw[arr] (apibranch) -- (sqlite.north);
\draw[arr] (apibranch) -| node[near start, xshift=3mm, lbl, above] {subprocess} (parser.north);
\draw[arr] (parser) -- (yandex);
\draw[arr] (api.east) -- (geocoder.west);
\end{tikzpicture}
\caption{Логическая архитектура программного комплекса}
\label{fig:architecture}
\end{figure}

\textbf{Режим разработки.} Клиент собирается Vite и обслуживается на порту~3000; запросы к \texttt{/api} проксируются на сервер (порт~3001), что даёт единый origin в браузере. В продуктивной сборке статика клиента отдаётся nginx из каталога \texttt{dist/}, API --- тем же Express-приложением за reverse-proxy (см.~раздел~3, подраздел о развёртывании).

\textbf{Разделение ответственности.} Парсер работает как изолированный Python-процесс: сервер только запускает стадии, читает файлы статуса и по завершении импортирует готовый CSV в JSON-хранилище. Визуализация целиком сосредоточена во frontend; сервер не генерирует HTML-страницы аналитики. Личный CSV аналитика хранится в браузере (\texttt{localStorage}) и не попадает на сервер --- это сознательное упрощение для роли «аналитик» без общего файлового хранилища.

\subsection{Проектирование потока данных и пайплайна парсера}

Сбор данных организован поэтапно --- в соответствии с теоретическим обоснованием поэтапного парсинга (раздел~1). Проектный поток приведён на рисунке~\ref{fig:parser-flow}.

\begin{figure}[htbp]
\centering
\begin{tikzpicture}[
    stage/.style={draw, rounded corners, minimum width=2.1cm, minimum height=0.75cm, align=center, font=\small},
    arr/.style={-{Stealth[length=2mm]}}
]
\node[stage] (links) {links};
\node[stage, right=0.55cm of links] (details) {details};
\node[stage, right=0.55cm of details] (etl) {etl};
\node[stage, right=0.55cm of etl] (qa) {qa};
\node[stage, right=0.55cm of qa] (ingest) {ingest};
\draw[arr] (links) -- (details);
\draw[arr] (details) -- (etl);
\draw[arr] (etl) -- (qa);
\draw[arr] (qa) -- (ingest);
\end{tikzpicture}
\caption{Основные стадии пайплайна сбора и подготовки данных}
\label{fig:parser-flow}
\end{figure}

\begin{enumerate}
    \item \textbf{links} --- обход поисковых сегментов города в браузере, накопление URL объявлений;
    \item \textbf{details} --- посещение карточек, извлечение цены, площадей, адреса, координат, параметров здания; поддержка буферизации и возобновления;
    \item \textbf{etl} --- нормализация адресов и координат, расчёт производных признаков (цена за м², расстояние до центра, возраст дома, one-hot кодирование категорий), отсечение заведомо некорректных записей;
    \item \textbf{qa} --- сверка с контрактом схемы: обязательные поля, дубликаты ссылок, выбросы по координатам;
    \item \textbf{ingest} --- загрузка итогового CSV в серверное JSON-хранилище с дедупликацией по идентификатору объявления.
\end{enumerate}

Промежуточные артефакты (списки ссылок, сырые и обработанные таблицы) сохраняются в каталоге данных парсера отдельно для каждого города. Сервер и UI опрашивают JSON-файл runtime-статуса: стадия, процент выполнения, признак капчи --- это позволяет администратору продолжить сбор вручную после прохождения проверки.

Контракт колонок датасета задаёт обязательный минимум (цена, комнаты, площадь, широта, долгота), идентификационные поля (ID, адрес, ссылка) и расширенный набор производных и категориальных признаков (тип дома, ремонт, лифты, парковка и др.). Единый контракт связывает выход ETL, проверку QA и ожидания дашборда.

\subsection{Проектирование модели данных}

Данные системы делятся на \emph{табличные снимки рынка} (основной объём) и \emph{служебные записи} в SQLite (компактные, транзакционные).

\subsubsection{Хранение объявлений}

Каждый город имеет каталог с двумя JSON-файлами:
\begin{itemize}
    \item \texttt{dataset.json} --- массив строк объявлений (\texttt{Record<string, unknown>});
    \item \texttt{dataset-meta.json} --- число строк, дата обновления, история загрузок (batchId, режим append/replace, статистика ingest).
\end{itemize}

Ограничения проектирования: не более 50~000 строк и 20~МБ на одну загрузку. При добавлении партии строки с одинаковым ID объявления сливаются: координаты и ссылки из разных выгрузок дополняют друг друга. Служебные поля (\texttt{\_\_uploadBatchId}, нормализованный тип дома) не экспортируются пользователю.

Ключевые группы столбцов датасета сведены в таблице~\ref{tab:dataset-fields}.

\begin{table}[htbp]
\centering
\small
\setlength{\tabcolsep}{3pt}
\begin{tabular}{|p{0.17\textwidth}|p{0.78\textwidth}|}
\hline
\textbf{Группа} & \textbf{Содержание} \\
\hline
Идентификация & ID объявления, ссылка на карточку, текстовый адрес \\
\hline
Цена и площадь & Цена (руб.), общая и жилая площадь, число комнат, цена за м² \\
\hline
Геопозиция & Широта, долгота, расстояние до центра города (км) \\
\hline
Здание & Год постройки, возраст дома, этаж, этажность, тип дома, лифты \\
\hline
Категории & One-hot признаки: ремонт, санузел, окна, парковка, двор, мебель, техника и др. \\
\hline
\end{tabular}
\caption{Логические группы полей датасета объявлений}
\label{tab:dataset-fields}
\end{table}

\subsubsection{Служебные таблицы SQLite}

В embedded-базе \texttt{app.db} хранятся сущности, для которых нужны транзакции и индексы (таблица~\ref{tab:sqlite}).

\begin{table}[htbp]
\centering
\small
\setlength{\tabcolsep}{3pt}
\begin{tabular}{|L{0.27\textwidth}|L{0.67\textwidth}|}
\hline
\textbf{Таблица} & \textbf{Назначение} \\
\hline
\texttt{users} & Учётные записи: логин, хеш пароля, роль (\texttt{observer}, \texttt{analyst}, \texttt{admin}, \texttt{manager}) \\
\hline
\texttt{parser\_jobs} & Задачи парсера: город, стадия, статус, PID процесса, пути к логу и конфигу, код завершения \\
\hline
\texttt{reverse\_geocode\_cache} & Кэш обратного геокодирования: ключ координат, адрес, срок действия (TTL, по умолчанию 14 суток) \\
\hline
\end{tabular}
\caption{Служебные таблицы SQLite}
\label{tab:sqlite}
\end{table}

Выбор JSON для массива объявлений и SQLite для пользователей и кэша обусловлен балансом простоты развёртывания и достаточной производительности при локальной работе: полный снимок города загружается в память клиента один раз (\texttt{bootstrap}), дальнейшая фильтрация идёт без обращений к серверу на каждый слайдер.

\subsection{Проектирование программного интерфейса сервера}

API построен в REST-стиле; авторизация --- JWT (Bearer-токен, срок жизни 12~часов). Эндпоинты сгруппированы по уровню доступа (таблица~\ref{tab:api}).

\begin{table}[htbp]
\centering
\footnotesize
\setlength{\tabcolsep}{3pt}
\textbf{Эндпоинты для авторизованных пользователей}\\[3pt]
\begin{tabular}{|L{0.34\textwidth}|L{0.11\textwidth}|L{0.49\textwidth}|}
\hline
\textbf{Путь} & \textbf{Метод} & \textbf{Назначение} \\
\hline
\texttt{/api/auth/login} & POST & Аутентификация, выдача JWT \\
\hline
\texttt{/api/me} & GET & Текущий пользователь \\
\hline
\texttt{/api/dataset/bootstrap} & GET & Полная выгрузка строк и сводки по городу \\
\hline
\texttt{/api/summary} & GET & Сводные показатели без всех строк \\
\hline
\texttt{/api/listings} & GET & Пагинация отфильтрованных объявлений \\
\hline
\texttt{/api/geocode/reverse} & GET & Обратное геокодирование одной точки \\
\hline
\end{tabular}

\medskip
\textbf{Эндпоинты analyst и admin}\\[3pt]
\begin{tabular}{|L{0.34\textwidth}|L{0.11\textwidth}|L{0.49\textwidth}|}
\hline
\textbf{Путь} & \textbf{Метод} & \textbf{Назначение} \\
\hline
\texttt{/api/export.csv} & GET & Экспорт CSV (analyst, admin) \\
\hline
\texttt{/api/admin/dataset} & POST/DEL & Загрузка и очистка серверного датасета \\
\hline
\texttt{/api/admin/dataset/}\newline\texttt{ingest-pipeline} & POST & Импорт CSV после ETL \\
\hline
\texttt{/api/admin/parser/jobs} & GET/POST & Список и запуск стадий парсера \\
\hline
\texttt{/api/admin/parser/jobs/}\newline\texttt{:id/resume} & POST & Продолжение после капчи \\
\hline
\end{tabular}

\medskip
\textbf{Эндпоинты manager}\\[3pt]
\begin{tabular}{|L{0.34\textwidth}|L{0.11\textwidth}|L{0.49\textwidth}|}
\hline
\textbf{Путь} & \textbf{Метод} & \textbf{Назначение} \\
\hline
\texttt{/api/manager/markets/overview} & GET & Сводка по всем городам \\
\hline
\texttt{/api/manager/users} & GET/POST & Управление учётными записями \\
\hline
\texttt{/api/manager/users/:id} & PATCH/DEL & Смена роли, пароля, удаление \\
\hline
\end{tabular}
\caption{Основные группы эндпоинтов REST API}
\label{tab:api}
\end{table}

Параметры фильтрации передаются в query string в едином формате, согласованном с клиентским модулем фильтров: диапазоны числовых полей, множественный выбор категорий, дополнительные условия по произвольным столбцам. Это гарантирует, что экспорт и постраничная выдача совпадают с тем, что пользователь видит на дашборде.

\subsubsection{Интеграция API Яндекса}

Для отображения карты объектов и обратного геокодирования координат в адрес система использует два сервиса экосистемы Яндекса: JavaScript API~2.1 на клиенте и API Геокодера на сервере (эндпоинт \texttt{/api/geocode/reverse}). Оба ключа выпускаются в едином кабинете разработчика; на рисунке~\ref{fig:yandex-api} показаны подключённые сервисы и переменные окружения \srce{VITE_YANDEX_MAPS_API_KEY} и \srce{YANDEX_GEOCODER_API_KEY}.

\begin{figure}[htbp]
\centering
\includegraphics[width=0.88\textwidth]{yandex-developer-console}
\caption{Кабинет разработчика Яндекса: ключи JavaScript API и API Геокодера}
\label{fig:yandex-api}
\end{figure}

На скриншоте видно, что для клиентской карты активирован JavaScript API, а для серверного эндпоинта обратного геокодирования --- API Геокодера; разделение ключей по уровням приложения снижает риск несанкционированного использования серверного ключа в браузере. Результаты геокодирования кэшируются в таблице \texttt{reverse\_geocode\_cache} (см.~табл.~\ref{tab:sqlite}).

\subsection{Проектирование ролевой модели}

Система рассчитана на несколько типов пользователей организации, изучающей рынок недвижимости. Матрица возможностей приведена в таблице~\ref{tab:roles}.

\begin{table}[htbp]
\centering
\small
\setlength{\tabcolsep}{3pt}
\begin{tabular}{|p{0.16\textwidth}|p{0.24\textwidth}|p{0.54\textwidth}|}
\hline
\textbf{Роль} & \textbf{Экран} & \textbf{Возможности} \\
\hline
observer & Дашборд & Серверный датасет; KPI, фильтры, фиксированные графики, карта; без таблицы, экспорта и своего CSV \\
\hline
analyst & Дашборд & Всё у observer плюс свой CSV, конструктор графиков, таблица, экспорт CSV/XLSX и PNG/PDF \\
\hline
admin & Дашборд + админ-панели & Управление серверным датасетом, запуск парсера, ingest; без блока «Обзор рынка» \\
\hline
manager & Отдельная консоль & Сводка по рынкам и CRUD пользователей; без доступа к дашборду \\
\hline
\end{tabular}
\caption{Роли пользователей и разграничение доступа}
\label{tab:roles}
\end{table}

Руководитель (\texttt{manager}) не может назначать роль \texttt{manager} другим пользователям --- это защита от эскалации прав. Наблюдатель получает упрощённый список объявлений вместо полной таблицы, чтобы снизить риск массового копирования данных (см.~рис.~\ref{fig:observer-dashboard} в подразделе~2.7.2). Администратор при пустом городе попадает на экран первичной настройки: предложение загрузить CSV или запустить сбор; панели ingest и парсера --- на рис.~\ref{fig:admin-server} и~\ref{fig:admin-parser}. Консоль руководителя --- на рис.~\ref{fig:manager-markets},~\ref{fig:manager-users}.

\subsection{Проектирование пользовательского интерфейса}

Интерфейс спроектирован по принципам, сформулированным в разделе~1: фильтры отделены от области результатов, обратная связь при изменении параметров --- немедленная, предусмотрены светлая и тёмная темы.

\subsubsection{Структура основного экрана}

Экран дашборда делится на шапку и рабочую зону. В шапке --- выбор города, переключатель источника данных (для аналитика), элементы темы, имя пользователя и выход.

Боковая панель фильтров на десктопе закреплена слева; на узком экране открывается поверх контента. На рис.~\ref{fig:dashboard-filters} показаны слайдеры диапазонов, чипы комнатности и типа дома, переключатели исключения этажей и конструктор дополнительных условий.

\screenshot[0.78]{dashboard-filters}{dashboard-filters}{Панель фильтров дашборда с активными параметрами отбора}

Первый блок основной колонки --- \textbf{обзор рынка}: медианы цены, цены за м², площадей, число объектов, краткие текстовые инсайты по активным фильтрам. Ниже располагается \textbf{сетка графиков} --- гистограммы, столбчатые диаграммы по категориям, точечные диаграммы с линией тренда; у аналитика и администратора доступен конструктор пользовательских графиков. Общий вид этих блоков для роли \texttt{analyst} с серверным датасетом Москвы --- на рис.~\ref{fig:dashboard-main}.

\screenshot{dashboard-main}{dashboard-main}{Основной экран дашборда (роль analyst, город Москва): шапка, KPI и сетка графиков}

Далее по вертикали следует \textbf{карта} --- кластеры объявлений на базе JavaScript API Яндекс.Карт, балуны с ценой и ссылкой, ограничение числа отображаемых точек для производительности (рис.~\ref{fig:dashboard-map}).

\screenshot{dashboard-map}{dashboard-map}{Карта объектов с кластерами и открытым балуном объявления}

Блок \textbf{сравнения сегментов} выводит два независимых набора фильтров и KPI рядом --- например, новостройки и вторичка (рис.~\ref{fig:dashboard-segments}).

\screenshot{dashboard-segments}{dashboard-segments}{Сравнение двух сегментов рынка (два набора фильтров и KPI)}

Завершает экран \textbf{таблица объявлений} для ролей analyst и admin: сортировка, настраиваемые колонки, геокодирование строк без адреса, экспорт (рис.~\ref{fig:dashboard-table}).

\screenshot[0.95]{dashboard-table}{dashboard-table}{Таблица отфильтрованных объявлений с настраиваемыми колонками (роль analyst)}

\subsubsection{Типовые сценарии}

\textbf{Аналитик:} вход $\rightarrow$ выбор серверного снимка или загрузка CSV $\rightarrow$ настройка фильтров (рис.~\ref{fig:dashboard-filters}) $\rightarrow$ просмотр медиан и графиков (рис.~\ref{fig:dashboard-main}) $\rightarrow$ проверка расположения на карте (рис.~\ref{fig:dashboard-map}) $\rightarrow$ экспорт таблицы или графика (рис.~\ref{fig:dashboard-table}).

\textbf{Наблюдатель:} вход $\rightarrow$ просмотр серверного датасета без операций записи и без полной таблицы. Интерфейс ограничен KPI, графиками и картой (рис.~\ref{fig:observer-dashboard}).

\screenshot{observer-dashboard}{observer-dashboard}{Дашборд роли observer: KPI, графики и карта без таблицы и экспорта}

\textbf{Администратор:} вход $\rightarrow$ панель серверных данных --- upload, ingest, удаление (рис.~\ref{fig:admin-server}) $\rightarrow$ запуск стадии парсера и мониторинг прогресса (рис.~\ref{fig:admin-parser}) $\rightarrow$ при капче --- ручное продолжение в браузере парсера $\rightarrow$ ingest готового CSV.

\screenshot{admin-server}{admin-server}{Админ-панель управления серверным датасетом (upload, ingest, метаданные города)}

\screenshot{admin-parser}{admin-parser}{Админ-панель запуска парсера и отображения прогресса задачи}

\textbf{Руководитель:} вход в консоль manager $\rightarrow$ вкладка «Рынки» --- объём данных, даты обновления, диапазон цен по городам (рис.~\ref{fig:manager-markets}) $\rightarrow$ вкладка «Пользователи» --- создание учёток, смена ролей (рис.~\ref{fig:manager-users}).

\screenshot{manager-markets}{manager-markets}{Консоль руководителя: вкладка «Рынки» со сводкой по городам}

\screenshot[0.88]{manager-users}{manager-users}{Консоль руководителя: вкладка «Пользователи» (создание учётных записей и ролей)}

\subsection{Обоснование технологического стека}

Выбор средств разработки согласован с требованиями раздела~1 и архитектурой на рис.~\ref{fig:architecture}. Сводка приведена в таблице~\ref{tab:stack}.

\begin{table}[htbp]
\centering
\small
\setlength{\tabcolsep}{3pt}
\begin{tabular}{|p{0.14\textwidth}|p{0.19\textwidth}|p{0.61\textwidth}|}
\hline
\textbf{Уровень} & \textbf{Технология} & \textbf{Обоснование} \\
\hline
Парсер & Python, Selenium~4, pandas & Динамический SPA Яндекс.Недвижимость; ETL в табличном виде~\cite{selenium,mckinney2023} \\
\hline
API & Node.js, Express~5, better-sqlite3 & Лёгкий REST-сервис, JWT, синхронный доступ к SQLite без отдельного сервера БД \\
\hline
Клиент & React~19, TypeScript, Vite~6 & SPA с типизацией и быстрой пересборкой~\cite{react,typescript,vite,sultanov2022} \\
\hline
Графики & Recharts~3 & Декларативные React-компоненты, синхронная перерисовка при смене фильтров~\cite{recharts} \\
\hline
Таблицы/CSV & PapaParse, SheetJS & Разбор и экспорт CSV на клиенте без нагрузки на сервер~\cite{papaparse} \\
\hline
Стили & Tailwind CSS & Утилитарная вёрстка, адаптивность, темы~\cite{tailwind} \\
\hline
Карта & Яндекс.Карты JS API~2.1 & Кластеризация и балуны; согласованность с экосистемой источника данных \\
\hline
Среда & VS Code, WebStorm & Мультиязычный проект (Python + TypeScript), единый репозиторий \\
\hline
\end{tabular}
\caption{Технологический стек и мотивировка выбора}
\label{tab:stack}
\end{table}

\textbf{Selenium вместо статического HTML-парсинга} --- карточки и выдача на \texttt{realty.yandex.ru} формируются JavaScript; без управляемого браузера выборка неполна (ср.~табл.~1.1). \textbf{SQLite вместо PostgreSQL} --- достаточно для учебного развёртывания и демонстрации ролевой модели; объём объявлений вынесен в JSON. \textbf{Фильтрация на клиенте} --- после bootstrap сервер не участвует в каждом изменении слайдера, что соответствует паттерну SPA~\cite{nesterov2024,filimonova2024}.

\subsection{Выводы по второму разделу}

Во втором разделе спроектирован программный комплекс для анализа рынка недвижимости на данных Яндекс.Недвижимость:
\begin{enumerate}
    \item сформулированы функциональные и нефункциональные требования, включая поэтапный сбор, мультгородность и ролевую модель;
    \item выделены четыре логических компонента (клиент, API, парсер, общий слой правил) и описаны потоки данных между ними;
    \item спроектированы модель хранения (JSON-снимки + SQLite), контракт полей датасета и группы REST-эндпоинтов;
    \item определены экраны и сценарии для ролей observer, analyst, admin и manager; показано подключение API Яндекс.Карт и Геокодера;
    \item обоснован технологический стек с опорой на теоретический анализ первого раздела.
\end{enumerate}
